\documentclass[11pt]{article}

\usepackage[T1]{fontenc}
\usepackage{lmodern}
\usepackage[margin=1in]{geometry}
\usepackage{microtype}
\usepackage{amsmath,amssymb,amsthm,mathtools}
\usepackage{enumitem}
\usepackage[hidelinks]{hyperref}

\newtheorem{theorem}{Theorem}[section]
\newtheorem{proposition}[theorem]{Proposition}
\newtheorem{lemma}[theorem]{Lemma}
\newtheorem{corollary}[theorem]{Corollary}
\theoremstyle{remark}
\newtheorem{remark}[theorem]{Remark}

\newcommand{\Cb}{C_b}
\newcommand{\B}{\mathcal B}
\newcommand{\Lat}{\operatorname{Lat}}
\newcommand{\clspan}{\overline{\operatorname{span}}}
\newcommand{\N}{\mathbb N}
\newcommand{\C}{\mathbb C}
\newcommand{\Q}{\mathbb Q}
\newcommand{\ellinf}{\ell_\infty}

\title{Operators Without Nontrivial Invariant Subspaces on $C_b(K)$\\
A Complete Classification and Literature Audit}
\author{}
\date{June 21, 2026}

\begin{document}
\maketitle

\begin{abstract}
We audit a proposed argument removing completeness from the final
$C(K)$ theorem of Cheng--Fang--Jiang.  The argument is correct: for every
metric space $K$ with infinite-dimensional $C_b(K)$, the space $C_b(K)$
admits a bounded operator without nontrivial closed invariant subspaces if
and only if $K$ is compact.  In fact, this follows from classical ingredients:
Conway's characterization of separable $C_b(K)$ for completely regular
spaces, Sobczyk's complementation theorem for $c_0$, and Read's construction
of operators without invariant subspaces.

We give a self-contained proof and a complete all-dimensional formulation.
For a nonempty topological space $K$, let $K_{\mathrm{cr}}$ be its
bounded-continuous-function reflection.  Then $C_b(K)$ supports an operator
with no nontrivial closed invariant subspaces if and only if
$K_{\mathrm{cr}}$ is compact metrizable and is either a singleton or an
infinite set.  Under the source paper's infinite-dimensional hypothesis, this
is equivalent to separability of $C_b(K)$.  For Tychonoff spaces, and hence
for metric spaces, the reflection is $K$ itself.  This result is separate from
Halmos's third problem, which it does not resolve.
\end{abstract}

\section{Context and verdict}

Let $K$ be a topological space and let $\Cb(K)$ denote the Banach space of
bounded continuous complex-valued functions on $K$, with the supremum norm.
For a Banach space $X$ and $T\in\B(X)$, write
\[
  \Lat(T)=\{M\subseteq X: M\text{ is closed and linear, and }T(M)\subseteq M\}.
\]
We call $T$ \emph{transitive} when $\Lat(T)=\{0,X\}$.

Cheng--Fang--Jiang proved that, for a complete metric space $K$ with
infinite-dimensional $\Cb(K)$, such a transitive operator exists on $\Cb(K)$
if and only if $K$ is compact \cite[Theorem 16.1]{CFJ}.  The proposed partial
argument removes completeness by combining a closed discrete sequence with
Tietze extension.

\paragraph{Verdict.}
The partial argument is mathematically correct.  The only point worth making
explicit is that Tietze extension is first applied to the real and imaginary
parts of a complex-valued function.  More importantly, the exact metric
statement is not an unresolved problem: it follows immediately from a
classical theorem of Conway that $\Cb(K)$ is separable exactly when a
Tychonoff space $K$ is compact metrizable \cite[Theorem V.6.6, p.~140]{Conway},
together with Read's and Sobczyk's theorems.  Sections~\ref{sec:audit}--
\ref{sec:classification} give a complete proof and a stronger classification.

\section{Audit of the proposed metric argument}\label{sec:audit}

The two lemmas in the proposed solution are both valid.

\begin{proposition}\label{prop:nonsep}
Let $X$ be a nonseparable Banach space.  Then every $T\in\B(X)$ has a
nontrivial closed invariant subspace.
\end{proposition}

\begin{proof}
Choose $0\neq x\in X$ and set
\[
  M_x=\clspan\{x,Tx,T^2x,\ldots\}.
\]
Then $M_x$ is nonzero, closed, and $T$-invariant.  It is separable because it
is generated by a countable set.  Since $X$ is nonseparable, $M_x\neq X$.
\end{proof}

The same observation gives a useful converse formulation.

\begin{corollary}\label{cor:transitive-separable}
If a Banach space supports a transitive bounded operator, then it is
separable.  Indeed, every nonzero vector is cyclic for that operator.
\end{corollary}

\begin{lemma}\label{lem:finite-dimensional}
Let $X$ be a nonzero finite-dimensional complex Banach space.  A transitive
operator exists on $X$ if and only if $\dim X=1$.
\end{lemma}

\begin{proof}
If $\dim X=1$, there are no nonzero proper linear subspaces.  If
$\dim X\geq2$, every $T\in\B(X)$ has an eigenvalue over $\C$.  The
corresponding eigenline is a nonzero proper closed invariant subspace.
\end{proof}

\begin{proposition}\label{prop:metric-noncompact}
If $K$ is a noncompact metric space, then $\Cb(K)$ is nonseparable.
\end{proposition}

\begin{proof}
Metric compactness is equivalent to sequential compactness.  Hence there is a
sequence in $K$ with no convergent subsequence.  After passing to a
subsequence, assume its terms $x_n$ are distinct and put
\[
  D=\{x_n:n\in\N\}.
\]
The set $D$ is closed and discrete in $K$: an accumulation point of $D$ would
produce a convergent subsequence.

Every bounded function on $D$ is continuous, so $\Cb(D)$ is isometric to
$\ellinf$.  The restriction map
\[
  R:\Cb(K)\longrightarrow \Cb(D),\qquad Rf=f|_D,
\]
is bounded and linear.  It is surjective.  Indeed, given
$a:D\to\C$ bounded, apply the real-valued Tietze extension theorem to
$\operatorname{Re}a$ and $\operatorname{Im}a$ separately; metric spaces are
normal and $D$ is closed.  Combining the two extensions gives a bounded
continuous complex-valued extension of $a$ to $K$.

Thus $\ellinf$ is a continuous linear image of $\Cb(K)$.  Since a continuous
image of a separable metric space is separable, while $\ellinf$ is
nonseparable, $\Cb(K)$ cannot be separable.
\end{proof}

\begin{corollary}\label{cor:metric-necessity}
Let $K$ be metric.  If $\Cb(K)$ supports a transitive bounded operator, then
$K$ is compact.
\end{corollary}

\begin{proof}
By Corollary~\ref{cor:transitive-separable}, $\Cb(K)$ is separable.  Proposition
\ref{prop:metric-noncompact} excludes noncompact $K$.
\end{proof}

This verifies the new part of the proposed proof.  Notice that completeness
plays no role.

\section{The classical separability theorem for $C_b(K)$}

A space is called \emph{Tychonoff} if it is completely regular and Hausdorff.
The following theorem is classical; the proof is included to make the logical
closure of the argument transparent.

\begin{theorem}[Conway]\label{thm:conway}
Let $K$ be a Tychonoff space.  Then $\Cb(K)$ is separable if and only if $K$ is
compact metrizable.
\end{theorem}

\begin{proof}
Suppose first that $\Cb(K)$ is separable, and choose a norm-dense sequence
$(f_n)_{n\geq 1}$ in $\Cb(K)$.  Define
\[
  \Phi:K\longrightarrow \C^{\N},\qquad
  \Phi(x)=(f_1(x),f_2(x),\ldots),
\]
where $\C^{\N}$ has the product topology.

We claim that $\Phi$ is a topological embedding.  If $x\neq y$, complete
regularity gives $g\in\Cb(K)$ with $g(x)=0$ and $g(y)=1$.  Choose $n$ with
$\lVert f_n-g\rVert_\infty<1/3$.  Then $f_n(x)\neq f_n(y)$, so $\Phi$ is
injective.

To see that the inverse on $\Phi(K)$ is continuous, let $x\in U$ with $U$ open
in $K$.  Complete regularity gives $g\in\Cb(K)$ such that
$g(x)=0$ and $g=1$ on $K\setminus U$.  Choose $n$ with
$\lVert f_n-g\rVert_\infty<1/4$.  Then
\[
  V=\{y\in K:|f_n(y)-f_n(x)|<1/2\}
\]
is an open neighborhood of $x$ contained in $U$.  Hence the coordinates
$f_n$ generate the topology of $K$, proving the claim.  Therefore $K$ is
metrizable.  If it were noncompact, Proposition~\ref{prop:metric-noncompact}
would imply that $\Cb(K)$ is nonseparable, a contradiction.  Thus $K$ is
compact metrizable.

Conversely, suppose $K$ is compact metrizable.  Choose a countable dense set
$\{x_j:j\geq1\}$ in $K$.  The unital algebra over $\Q+i\Q$ generated by the
real-valued functions $d(\,\cdot\,,x_j)$ is countable, self-adjoint, and
separates points.  The complex Stone--Weierstrass theorem therefore makes it
uniformly dense in $C(K)=\Cb(K)$.  Thus $\Cb(K)$ is separable.
\end{proof}

For metric spaces, Theorem~\ref{thm:conway} directly says that the
completeness hypothesis in \cite[Theorem 16.1]{CFJ} is unnecessary.  The
Tietze proof in Section~\ref{sec:audit} is an elementary proof of the only
nontrivial direction needed for that corollary.

\section{The compact-metric existence direction}\label{sec:existence}

For completeness, we now supply the operator-theoretic direction rather than
merely invoking Cheng--Fang--Jiang.

\begin{proposition}\label{prop:compact-existence}
Let $L$ be an infinite compact metric space.  Then $C(L)$ admits a bounded
operator with no nontrivial closed invariant subspaces.
\end{proposition}

\begin{proof}
First we exhibit an isometric copy of $c_0$ in $C(L)$.  Since $L$ is infinite
and compact metric, it contains a sequence of distinct points converging to a
point $p\in L$.  Passing to a subsequence, write the points as $x_n$ and set
$r_n=d(x_n,p)>0$ so that
\[
  r_{n+1}<\frac14 r_n\qquad(n\geq1).
\]
Define
\[
  u_n(x)=\max\left\{0,1-\frac{4d(x,x_n)}{r_n}\right\},\qquad x\in L.
\]
Then $0\leq u_n\leq1$, $u_n(x_n)=1$, and the supports of the $u_n$ are
pairwise disjoint.  Indeed, for $m>n$,
\[
  d(x_n,x_m)\geq r_n-r_m>\frac34r_n,
\]
whereas points in both supports would force
$d(x_n,x_m)\leq r_n/4+r_m/4<5r_n/16$.

For $a=(a_n)\in c_0$, set
\[
  J(a)=\sum_{n=1}^{\infty}a_nu_n.
\]
The series converges uniformly because the supports are disjoint and
$a_n\to0$.  Moreover
\[
  \lVert J(a)\rVert_\infty=\sup_n|a_n|,
\]
so $J$ is an isometric embedding of $c_0$ into $C(L)$.

The space $C(L)$ is separable.  Sobczyk's theorem says that every copy of
$c_0$ in a separable Banach space is complemented \cite{Sobczyk}.  Hence
\[
  C(L)\cong c_0\oplus Y
\]
for some separable Banach space $Y$.  Read proved that every separable complex
Banach space of the form $c_0\oplus Y$ admits a bounded operator without
nontrivial closed invariant subspaces \cite{Read}.  Transporting that operator
through the above isomorphism gives the required operator on $C(L)$.
\end{proof}

\begin{remark}
The complemented form of Read's theorem used here is exactly the
Read--Sobczyk route invoked by Cheng--Fang--Jiang.  No explicit formula for
Read's operator is needed for the classification.
\end{remark}

\section{A full classification for arbitrary topological spaces}
\label{sec:classification}

The metric theorem has a natural maximal extension.  For an arbitrary
topological space $K$, define the evaluation map
\[
  e_K:K\longrightarrow \C^{\Cb(K)},\qquad
  e_K(x)=(f(x))_{f\in\Cb(K)},
\]
and let
\[
  K_{\mathrm{cr}}=e_K(K)
\]
with the subspace topology inherited from the product.  The space
$K_{\mathrm{cr}}$ is Tychonoff.  It is the reflection of $K$ detected by its
bounded continuous scalar-valued functions.

\begin{lemma}\label{lem:reflection}
Pullback along $e_K$ is an isometric isomorphism
\[
  U_K:\Cb(K_{\mathrm{cr}})\longrightarrow \Cb(K),
  \qquad U_Kg=g\circ e_K.
\]
\end{lemma}

\begin{proof}
Because $e_K$ is onto $K_{\mathrm{cr}}$, pullback is isometric and injective.
For each $f\in\Cb(K)$, the corresponding coordinate projection
\[
  \widehat f:K_{\mathrm{cr}}\to\C,
  \qquad \widehat f\bigl((z_h)_{h\in\Cb(K)}\bigr)=z_f,
\]
is bounded and continuous, and $\widehat f\circ e_K=f$.  Thus $U_K$ is
surjective.
\end{proof}

\begin{theorem}[Complete classification]\label{thm:main}
Let $K$ be a nonempty topological space.  There is an operator
$T\in\B(\Cb(K))$ with $\Lat(T)=\{0,\Cb(K)\}$ if and only if the
bounded-function reflection $K_{\mathrm{cr}}$ is compact metrizable and
$K_{\mathrm{cr}}$ is either a singleton or infinite.
\end{theorem}

\begin{proof}
Suppose first that such an operator exists.  Corollary
\ref{cor:transitive-separable} shows that $\Cb(K)$ is separable.  By
Lemma~\ref{lem:reflection}, $\Cb(K_{\mathrm{cr}})$ is separable.  Since
$K_{\mathrm{cr}}$ is Tychonoff, Theorem~\ref{thm:conway} implies that
$K_{\mathrm{cr}}$ is compact metrizable.  If $K_{\mathrm{cr}}$ is finite of
cardinality $m$, then
\[
  \Cb(K)\cong C(K_{\mathrm{cr}})\cong\C^m.
\]
Lemma~\ref{lem:finite-dimensional} forces $m=1$.

Conversely, if $K_{\mathrm{cr}}$ is a singleton, then
$\Cb(K)\cong\C$, so there is no nonzero proper subspace to be invariant.  If
$K_{\mathrm{cr}}$ is infinite compact metrizable, Lemma
\ref{lem:reflection} identifies $\Cb(K)$ isometrically with
$C(K_{\mathrm{cr}})$, and Proposition~\ref{prop:compact-existence} supplies
a transitive operator on the latter.  Conjugating by the isometric
isomorphism gives one on $\Cb(K)$.
\end{proof}

\begin{corollary}[Infinite-dimensional form]\label{cor:infinite}
Let $K$ be a topological space and suppose $\Cb(K)$ is infinite-dimensional.
The following are equivalent:
\begin{enumerate}[label=\textup{(\roman*)},leftmargin=2.2em]
  \item $\Cb(K)$ supports a bounded operator without nontrivial closed
        invariant subspaces;
  \item $\Cb(K)$ is separable;
  \item $K_{\mathrm{cr}}$ is compact metrizable.
\end{enumerate}
\end{corollary}

\begin{proof}
The equivalence of \textup{(ii)} and \textup{(iii)} follows from
Lemma~\ref{lem:reflection} and Theorem~\ref{thm:conway}.  Since
$\Cb(K)$ is infinite-dimensional, $K_{\mathrm{cr}}$ is infinite, so
Theorem~\ref{thm:main} gives \textup{(i)$\Leftrightarrow$(iii)}.
\end{proof}

\begin{corollary}[Tychonoff version]\label{cor:tych}
Let $K$ be a nonempty Tychonoff space.  Then $\Cb(K)$ admits a bounded
operator without nontrivial closed invariant subspaces if and only if $K$ is
compact metrizable and $K$ is either a singleton or infinite.
\end{corollary}

\begin{proof}
For Tychonoff $K$, the evaluation map $e_K$ is a homeomorphism from $K$ onto
$K_{\mathrm{cr}}$.  Apply Theorem~\ref{thm:main}.
\end{proof}

\begin{corollary}[Metric version, all dimensions]\label{cor:metric-all}
Let $K$ be a nonempty metric space.  Then $\Cb(K)$ admits a bounded operator
without nontrivial closed invariant subspaces if and only if $K$ is compact
and $K$ is either a singleton or infinite.
\end{corollary}

\begin{corollary}[Metric version under the source hypothesis]\label{cor:metric}
Let $K$ be an arbitrary metric space and suppose $\Cb(K)$ is
infinite-dimensional.  Then $\Cb(K)$ admits a bounded operator without
nontrivial closed invariant subspaces if and only if $K$ is compact.
\end{corollary}

This is precisely the desired removal of completeness.

\begin{remark}[Why a separation hypothesis matters]
The conclusion ``$K$ must be compact'' cannot be extended verbatim to all
topological spaces.  Let $P$ be an infinite set with the particular-point
topology: fix $p\in P$ and declare every nonempty open set to contain $p$.
Then $P$ is noncompact, but every continuous map from $P$ to a Hausdorff space
is constant.  If
\[
  K=[0,1]\sqcup P
\]
is the topological disjoint union, then $K$ is noncompact while
\[
  \Cb(K)\cong C([0,1])\oplus_\infty\C
       \cong C([0,1]\sqcup\{*\}).
\]
The latter is a separable infinite-dimensional $C(L)$ space with $L$ compact
metric, so it supports a transitive operator.  The correct arbitrary-space
criterion is therefore Theorem~\ref{thm:main}, phrased in terms of
$K_{\mathrm{cr}}$.
\end{remark}

\section{Literature outcome and scope}

The literature search leads to the following conclusions.

\begin{enumerate}[leftmargin=2.2em]
  \item The proposed Tietze argument is correct and supplies an elementary
        proof for arbitrary metric spaces.
  \item A full answer was already implicit in classical literature.  Conway's
        theorem identifies separability of $\Cb(K)$ with compact metrizability
        for Tychonoff $K$; Read and Sobczyk supply the transitive operator on
        the compact-metric side.  Thus there is no counterexample to the
        proposed metric statement.
  \item The reflection theorem above packages the same ingredients for
        arbitrary topological spaces and also records the finite-dimensional
        exceptions.  We have not located an earlier source that states this
        exact packaging, but its ingredients are classical, so no novelty claim
        should be made without a more exhaustive bibliographic investigation.
  \item None of these statements settles Halmos's third problem for bounded
        invertible operators.  Recent work still treats that problem through
        additional spectral hypotheses; see \cite{JFJ}.
\end{enumerate}

\begin{thebibliography}{99}

\bibitem{CFJ}
L. Cheng, J. Fang, and C. Jiang,
\emph{On Halmos' third problem on Banach spaces},
Sci. China Math. \textbf{68} (2025), 399--420.
\href{https://doi.org/10.1007/s11425-023-2239-7}
{doi:10.1007/s11425-023-2239-7}.

\bibitem{Conway}
J.~B. Conway,
\emph{A Course in Functional Analysis}, 2nd ed.,
Graduate Texts in Mathematics 96, Springer, New York, 1990.

\bibitem{JFJ}
C. Jiang, J. Fang, and K. Ji,
\emph{Cowen--Douglas operators and the third of Halmos' ten problems},
Israel J. Math. \textbf{265} (2025), 379--427.
\href{https://doi.org/10.1007/s11856-024-2669-y}
{doi:10.1007/s11856-024-2669-y}.

\bibitem{Read}
C.~J. Read,
\emph{The invariant subspace problem on some Banach spaces with separable
 dual},
Proc. Lond. Math. Soc. (3) \textbf{58} (1989), no.~3, 583--607.
\href{https://doi.org/10.1112/plms/s3-58.3.583}
{doi:10.1112/plms/s3-58.3.583}.

\bibitem{Sobczyk}
A. Sobczyk,
\emph{Projection of the space $m$ on its subspace $c_0$},
Bull. Amer. Math. Soc. \textbf{47} (1941), 938--947.

\bibitem{Tietze}
H. Tietze,
\emph{Ueber Funktionen, die auf einer abgeschlossenen Menge stetig sind},
J. Reine Angew. Math. \textbf{145} (1915), 9--14.

\end{thebibliography}

\end{document}
